{\large\textbf{Diffraction of X-rays by Structured and Evolving Targets (10 points)}}

\textbf{Note}\\
1. Vectors are denoted by bold symbols (e.g., $\mathbf{r}$, $\mathbf{q}$ ).\\
2. Assume that absorption is neglected and that the electric field is polarized
perpendicular to the plane of incidence.

X-ray diffraction patterns arise because many tiny ``wave sources'' within a
crystal interfere, and we can predict this by adding their complex amplitudes
with the correct phases. Consider a monochromatic wave characterized by a
(real) amplitude $A \geq 0$ and a phase $\phi$. We define the
\textbf{complex amplitude} $\tilde{A}$ as
\[
\tilde{A} = Ae^{i\phi},
\]
so that the (real) amplitude of the wave is the magnitude (absolute value) of
$\tilde{A}$, and $\phi$ is its phase. Thus $\tilde{A}$ conveniently encodes both
magnitude and phase in a single complex number. In this problem set, we define
a dimensionless intensity by the squared magnitude of the total complex
amplitude:
\[
I = |\tilde{A}|^2 = A^2.
\]
Common experimental and geometric proportionality factors, such as detector
response, incident-beam normalization, and common propagation factors, are
absorbed into this definition. By contrast, the single-electron scattering
amplitude $f_0$ is retained explicitly as the amplitude scale for scattering
from one point electron.

When a crystalline material is exposed to an incident wave, the wave is
diffracted by the crystal lattice and the diffracted parts interfere with one
another. The intensity of the resulting diffracted wave can be calculated by
adding the complex amplitudes of the individual diffracted waves, taking into
account the phase differences between them, and then computing the squared
magnitude of the resulting total complex amplitude. Diffraction primarily
arises from interactions with electrons, and contributions from heavier
particles such as nuclei are typically negligible. The amplitude of the wave
diffracted by a single point electron depends only on $R = |\mathbf{R}|$, the
distance from the electron to the detector. Since $R$ is much larger than the
dimensions of the sample, its variation across the sample can be neglected.
Therefore, the total diffracted wave can be accurately determined by properly
accounting for the phase differences between individual diffracted waves, while
their amplitudes are assumed to be constant.

Let $\mathbf{k}_i$ and $\mathbf{k}_f$ denote the wavevectors of the incident and
diffracted waves, respectively. An incident plane wave with wavevector
$\mathbf{k}_i$ is diffracted into a wave with wavevector $\mathbf{k}_f$. The
\textbf{momentum transfer} is defined as
\[
\mathbf{q} = \mathbf{k}_f - \mathbf{k}_i
\]
and their magnitudes are assumed to be equal, because the wavelength is
unchanged:
\[
k \equiv |\mathbf{k}_i| = |\mathbf{k}_f| = \frac{2\pi}{\lambda}.
\]

\textbf{Part A: Diffraction from two electrons treated as point particles (2.0 pts)}

Consider two point electrons located at positions $\mathbf{r}_1$ and
$\mathbf{r}_2$, and define $\mathbf{r} \equiv \mathbf{r}_2 - \mathbf{r}_1$. A
detector is located at $\mathbf{P}$, and we define
$\mathbf{R} \equiv \mathbf{P} - \mathbf{r}_1$. A plane wave
$E_i(\mathbf{r}) \propto e^{i\mathbf{k}_i\cdot\mathbf{r}}$ is incident on the two
electrons, and the diffracted wave is observed in the far field along the
direction $\mathbf{k}_f$. In part \textbf{A}, the common time-dependent factor
$e^{-i\omega t}$ is suppressed, since only relative spatial phases are relevant.

\textbf{A.1}\quad Under the far-field approximation,
$R \equiv |\mathbf{R}| \gg |\mathbf{r}|$, keep only the leading-order term in
$|\mathbf{r}|/R$ and write the outgoing geometric path difference,
$\Delta L_\mathrm{out} \equiv |\mathbf{P} - \mathbf{r}_1| - |\mathbf{P} - \mathbf{r}_2|$,
in terms of $\mathbf{r}$ and $\mathbf{k}_f$, or equivalently
$\mathbf{k}_f/k_f$. \hfill 0.6 pt

\textbf{A.2}\quad Using your result from A.1 and accounting for the
position-dependent phase of the incident wave at $\mathbf{r}_1$ and
$\mathbf{r}_2$, find the phase difference between the two diffracted
contributions at the detector, expressed in terms of $\mathbf{q}$ and
$\mathbf{r}$. The phase difference is defined as
$\Delta\phi \equiv \phi_1 - \phi_2$, where $\phi_1$ and $\phi_2$ are the phases
of the contributions from the electrons at $r_1$ and $r_2$. \hfill 0.4 pt

\textbf{A.3}\quad Express the total complex amplitude of the diffracted wave
from these two electrons in terms of $\mathbf{q}$ and $\mathbf{r}$. You may
ignore any overall common phase factor, since it does not affect the intensity.
Assume that the real amplitude of a diffracted wave from a single point
electron is a constant $f_0$, independent of position. \hfill 0.6 pt

\textbf{A.4}\quad Express the intensity of diffracted waves from those two
electrons in terms of $\mathbf{q}$ and $\mathbf{r}$. \hfill 0.4 pt

\textbf{Part B: Finite longitudinal coherence (phase-jump model) (2.3 pts)}

Consider two point-like electrons located at positions $\mathbf{r}_1$ and
$\mathbf{r}_2$, with $\mathbf{r} \equiv \mathbf{r}_2 - \mathbf{r}_1$. A beam of
wavelength $\lambda_0$ (so $k = 2\pi/\lambda_0$ and $\omega = 2\pi c/\lambda_0$)
illuminates the electrons, and the diffracted wave is observed in the far field
along $\mathbf{k}_f$. We model the incident field as a plane wave with a
time-dependent random phase,
\[
E_i(\mathbf{r}, t) = A\ \exp[i\,(\mathbf{k}_i\cdot\mathbf{r} - \omega t + \phi(t))]\,,
\]
where $\phi(t)$ is \textit{piecewise constant} and undergoes
\textit{random phase jumps} at regular time intervals of duration
\[
t_0 \equiv \frac{L_0}{c},
\]
with $L_0$ the (given) longitudinal coherence length. At the beginning of each
interval of length $t_0$, $\phi(t)$ is reset to a new independent value
uniformly distributed on $[0, 2\pi)$. Let $I_0$ denote the (time-averaged)
intensity at the detector that would be obtained from a \textit{single} electron
in the same geometry. In this problem, the detector is assumed to measure a
time-averaged intensity. That is, it does not resolve the individual random
phase jumps. Instead, it records the average of the instantaneous intensity
over a time much longer than the phase-jump interval $t_0$:
$\langle I\rangle_t \equiv \left\langle |E(t)|^2 \right\rangle_t$ . Here
$\langle\cdots\rangle_t$ denotes an average over many phase-jump intervals.

\textbf{B.1}\quad Using the phase-jump model above, derive the
\textit{time-averaged} total intensity $\langle I\rangle_t$ at the detector from
the two electrons. Your final result should be written in terms of $I_0$,
$\mathbf{q} \equiv \mathbf{k}_f - \mathbf{k}_i$, the separation vector
$\mathbf{r}$, the wavelength $\lambda_0$, and the coherence length $L_0$.
Assume the detector averages over times much longer than $t_0$. \hfill 2.3 pt

\textbf{Part C: Non-point particle effect (1.0 pts)}

An electron is often treated as a classical point particle, but in a more
realistic model its charge may be regarded as being distributed over a finite
spatial region. Consider two idealized charge distributions: (i) an ideal point
charge located at $\mathbf{r} = 0$, whose scattering amplitude is
$A_1(\mathbf{q}) = Q_0$, (ii) an extended Gaussian charge distribution
$\rho_2(\mathbf{r}) = \rho_0 \exp\left(-\frac{r^2}{R_0^2}\right), \qquad r = |\mathbf{r}|.$

The constants $Q_0$ and $\rho_0$ are chosen so that the total charge is the same
in both cases:
\[
Q_0 = \int \rho_2(\mathbf{r})\, d^3r.
\]
Here $d^3r$ denotes the volume element in three-dimensional space. In Cartesian
coordinates, $d^3r = dxdydz$ and $\int_{\mathbb{R}^3}$ means integration over all
space. Useful identities (may be used without proof):
\[
\int_0^\infty e^{-r^2/R_0^2}\, 4\pi r^2\, dr = \pi^{3/2}R_0^3, \int_{\mathbb{R}^3} e^{-\alpha r^2} e^{i\mathbf{k}\cdot\mathbf{r}}\, d^3\mathbf{r} = \left(\frac{\pi}{\alpha}\right)^{3/2} \exp\left(-\frac{k^2}{4\alpha}\right), \qquad \alpha > 0.
\]

\textbf{C.1}\quad Obtain the relation among $Q_0$, $\rho_0$, and $R_0$.
\hfill 0.4 pt

\textbf{C.2}\quad Evaluate the amplitude \hfill 0.4 pt
\[
A_2(\mathbf{q}) \equiv \int_{\mathbb{R}^3} \rho_2(\mathbf{r})e^{i\mathbf{q}\cdot\mathbf{r}}\, d^3\mathbf{r}
\]
and compare it with the point-charge amplitude $A_1(\mathbf{q}) = Q_0$.

\textbf{C.3}\quad Estimate the ratio of the diffracted intensities,
$\dfrac{I_2}{I_1}$, for these two idealized cases when $q = \dfrac{2}{R_0}$.
\hfill 0.2 pt

\textbf{Part D: Diffraction from a film with non-flat surface morphology (2.4 pts)}

Imagine that the surface of a film (i.e., the top atomic layers) is not
perfectly flat, but exhibits surface roughness. A common model is to assume
that the local film thickness (measured in monolayers) follows a Gaussian
distribution. Let $N$ denote the local number of completed monolayers within a
lateral coherence area of the beam. We assume that $N$ varies across the
surface and is normally distributed with mean $\bar{N}$ and standard deviation
$\sigma$ (both in units of monolayers):
\[
P(N) = \frac{1}{\sqrt{2\pi}\sigma} \exp\left[-\frac{(N - \bar{N})^2}{2\sigma^2}\right].
\]
(For the purpose of evaluating averages, you may treat $N$ as a continuous
variable.) Let $d$ be the spacing between adjacent atomic layers (monolayers),
and let $q_z$ denote the component of the scattering vector
$\mathbf{q} = \mathbf{k}_f - \mathbf{k}_i$ normal to the flat surface of the
film, i.e., $q_z = \mathbf{q}\cdot\hat{\mathbf{z}}$. For an integer number of
monolayers $N$, the scattering amplitude is
\[
A_N(q_z) = \sum_{n=0}^{N-1} e^{iq_z nd} = \frac{1 - e^{iq_z Nd}}{1 - e^{iq_z d}}.
\]
When averaging over the Gaussian thickness distribution, we treat $N$ as a
continuous variable and use the closed-form expression
\[
A_N(q_z) \equiv \frac{1 - e^{iq_z Nd}}{1 - e^{iq_z d}}
\]
as the corresponding continuous extension. Assume the measured diffraction
intensity is obtained from the \textit{coherent} average amplitude,
\[
I(q_z) \equiv \left|\langle A(q_z)\rangle\right|^2, \qquad \langle A(q_z)\rangle = \int_{-\infty}^{\infty} P(N)\, A_N(q_z)\, dN.
\]

\textbf{D.1}\quad Calculate the intensity ratios \hfill 2.4 pt
\[
\frac{I(q_z, \sigma = 0.4, \bar{N} = 5)}{I(q_z, \sigma = 0, \bar{N} = 5)}
\]
at
\[
q_z = \frac{\pi}{2d} \qquad \mathrm{and} \qquad q_z = \frac{2\pi}{d},
\]
respectively. If needed, evaluate the second case by taking the appropriate
limit.

\textbf{Part E: Diffraction from a film with evolving surface morphology (2.3 pts)}

Imagine a thin film with a simple cubic structure grown on a substrate in a
layer-by-layer mode, \textit{i.e.}, each monolayer is completed before the next
monolayer starts to grow. Let $d$ be the spacing between adjacent atomic layers
(monolayers), and let $q_z$ denote the component of
$\mathbf{q} = \mathbf{k}_f - \mathbf{k}_i$ normal to the flat surface of the
film, \textit{i.e.,} $q_z = \mathbf{q}\cdot\hat{\mathbf{z}}$. As the film is
being grown, the thickness of the film changes, and so does the diffraction
intensity at $\mathbf{q} = \frac{\pi}{d}\hat{\mathbf{z}}$. The calculation is
performed at the specified out-of-plane momentum transfer, and all monolayer
contributions are coherently summed.

\textbf{E.1}\quad If the film starts to grow at $t = 0$ and the time required to
complete one monolayer is $t_0$, obtain the diffraction intensity ratio
\hfill 2.3 pt
\[
\frac{I(t = 0.8t_0)}{I(t = 3.6t_0)}
\]
measured at $\mathbf{q} = \frac{\pi}{d}\hat{\mathbf{z}}$. Assume that during
each monolayer-growth interval, the fractional coverage of the top layer
increases linearly from 0 to 1, so that at time $t = (N + \theta)t_0$, there are
$N$ completed monolayers and a fractional coverage $\theta$ of the next
monolayer.
